\documentclass[a4paper,11pt]{article}

% ---------------------------------------------------------------------
% openPDKcreator -- technical report. Plain article class (this is a
% software-tool writeup grounded in one real, downloaded PDK, not a
% novel research result), A4, single file for easy Overleaf import.
% Do not compile locally -- imported into Overleaf and compiled there.
% ---------------------------------------------------------------------

\usepackage[margin=2.5cm]{geometry}
\usepackage{amsmath,amssymb}
\usepackage{booktabs}
\usepackage{longtable}
\usepackage{listings}
\usepackage{xcolor}
\usepackage{hyperref}
\usepackage{enumitem}

\lstset{
  basicstyle=\ttfamily\small,
  breaklines=true,
  frame=single,
  columns=fullflexible,
  backgroundcolor=\color{gray!8},
}

\title{openPDKcreator: A General-Purpose Toolkit for Reading, Editing,\\
and Round-Trip Exporting Process Design Kits}
\author{Rodrigo Picos\\
\small Electronic Technology, Department of Industrial Engineering and Construction,\\
\small Universitat de les Illes Balears, Palma, Spain\\
\small \href{mailto:rpicosg@gmail.com}{rpicosg@gmail.com}}
\date{\today}

\begin{document}
\maketitle

\begin{abstract}
Process Design Kits (PDKs) bundle a foundry's technology data --
layer definitions, design rules, device models, standard-cell views,
and tool-specific configuration -- into a large, heterogeneous tree of
proprietary and semi-standard file formats. Existing open-source PDK
tooling either targets one specific tool's native format in isolation
(a DRC deck, a LEF file, a Liberty library) or treats the PDK as an
opaque blob to be consumed, not edited. \textbf{openPDKcreator} is a
general-purpose toolkit that reads, cross-references, edits, and
losslessly re-exports a real, complete, open-source PDK --
IHP's SG13G2 130\,nm process, released under the open\_pdks
convention -- across all of its native tool formats at once: KLayout
layer properties, Magic technology files, LEF, DRC decks, Liberty
timing libraries, GDS structural summaries, CDL/SPICE/Verilog
netlists, ngspice model cards, and xschem symbols. Every editable
domain has a verified, byte-identical write-back path, checked both
per-file and via a whole-PDK, two-generation round-trip export
covering all 5{,}345 real files (744\,MB) with zero structural or
byte-level drift. The toolkit also supports linking user-authored
Verilog/Verilog-A behavioral models into the same cell-centric view
that aggregates a cell's native LEF/CDL/SPICE/Verilog/Liberty/GDS
representations. This report describes the toolkit's architecture,
the real-data verification discipline used throughout its
development, its current coverage across each native format, and the
work still open toward the stated end goal of authoring PDK content
from scratch, not just editing an existing one.
\end{abstract}

\tableofcontents

\section{Introduction}

A PDK is not one file format but a federation of them: a foundry
ships design-rule decks for whichever verification tool it targets
(Magic, KLayout, Calibre, \ldots), device models for whichever
simulator its customers use (ngspice, Xyce, \ldots), standard-cell
views in whatever a digital flow expects (LEF/Liberty/Verilog), and
schematic symbols for whatever schematic capture tool a designer
prefers (xschem, Qucs-S, \ldots). Each of these formats predates the
open-source EDA ecosystem that now consumes them, and each was
designed around one tool's own needs, not a shared, general data
model. The practical consequence is that most PDK tooling is
narrowly scoped: a LEF parser reads LEF, a DRC-deck linter reads one
tool's DRC syntax, and cross-referencing "what does this one cell
look like across every one of these formats" is left to the human
designer, manually opening each file in its own tool.

\textbf{openPDKcreator} takes a different approach: one project, one
data model per format, one cell-centric aggregation view across all
of them, grounded end-to-end in a real, complete, publicly available
PDK rather than synthetic or partial test fixtures. This report
documents that project as it stands: what real formats it reads,
what it can edit and losslessly write back, how that write-back
fidelity was verified at full-PDK scale, and the still-open work
toward the stated longer-term goal -- generating PDK content from
scratch, not only reading and editing an existing one.

\subsection{Why a real PDK, not fixtures}

Every parser and writer in this project was built and verified
directly against IHP's own downloaded SG13G2 data -- never a hand-written
sample file standing in for the real format. This discipline
surfaced real structural facts a synthetic fixture would not have:
that a single Magic technology type's cross-section recipe can be
split across multiple, textually distant blocks sharing one name
(Section~\ref{sec:magic}); that some writers must regenerate a
declaration line rather than patch it in place, because the original
can wrap across multi-line continuations in ways that are unsafe to
token-patch (Section~\ref{sec:netlist-writeback}); and that a naive
trailing-newline normalization can silently mask a real byte-level
bug in exactly the files whose test coverage most needed to catch it
(Section~\ref{sec:verification}).

\subsection{Relationship to open\_pdks}

IHP's own README states plainly that its PDK ``already comes
pre-built'' in the format described by the \texttt{open\_pdks}
project (\url{https://github.com/RTimothyEdwards/open_pdks}) --
direct, authoritative confirmation that this project's target
directory layout
(\texttt{\{libs.doc,libs.qa,libs.ref,libs.tech\}}, with
\texttt{libs.ref/<family>/\{cdl,gds,lef,lib,mag,spice,verilog,\allowbreak
xschem\}} and \texttt{libs.tech/<tool>/}) is the real, standard
convention, not a project-specific guess. One concrete divergence was
found and is documented rather than silently normalized away: the
canonical \texttt{open\_pdks} example library name carries an
explicit foundry tag segment (\texttt{<family>\_fd\_<type>}), while
IHP's real library names omit it
(\texttt{sg13g2\_pr}/\texttt{sg13g2\_io}/\texttt{sg13g2\_sram}/\allowbreak
\texttt{sg13g2\_stdcell}).

\section{Architecture}

\subsection{Directory layout and data flow}

The tool is organized around three layers: \texttt{pdklib/} (one module
per native format, each exposing a real, typed parse of that
format's structured content), \texttt{gui/} (a Tk application whose
tabs mirror the \texttt{pdklib/} module boundaries, grouped by what they
represent rather than left flat), and \texttt{export.py}/\texttt{main.py}
(CLI commands and a write-back export path that reads the same
parsed, possibly-edited in-memory objects the GUI edits and
serializes them back to their native format).

A downloaded copy of the real PDK lives under \texttt{data/}, fetched
via a depth-1, blob-filtered, sparse \texttt{git clone} of
IHP's public repository that deliberately excludes its own git
submodules (separate third-party repositories for digital synthesis
and EM solvers, unrelated to PDK structure). This directory is
gitignored -- the toolkit's own code is what is version-controlled,
not a redistribution of IHP's data.

\subsection{Cell-centric aggregation}

The \textbf{By Cell} view (\texttt{pdklib/cells.py}, \texttt{gui/cell\_hub\_view.py})
is the toolkit's central cross-reference: for one real cell name, it
aggregates whichever of LEF, CDL, SPICE, Verilog, Liberty, and GDS
representations exist for it, scanning every real file under each
view directory rather than assuming one fixed file-per-cell
convention -- a necessary generalization, since IHP's own families
use two different real organizations for the same view kind
(\texttt{sg13g2\_io}/\texttt{sg13g2\_stdcell} pack every cell for a
view into one combined file; \texttt{sg13g2\_sram} ships one real
file per hard macro). All parsing runs through one shared,
application-owned cache, so an edit made from one tab (e.g.\ a LEF
pin edited via \textbf{By Cell}) is immediately visible, as the exact
same in-memory object, from every other tab that shows the same
data -- confirmed by object-identity checks, not just value equality,
in the project's own regression tests.

\section{Native format coverage}

\subsection{KLayout layers}

\texttt{pdklib/layers.py} parses KLayout's \texttt{.lyp} layer-properties
XML via a hand-rolled, nesting-depth-aware text scanner (rather than
\texttt{ElementTree}) specifically so that each real layer's own
source line range can be tracked -- a prerequisite for the surgical,
line-range write-back in \texttt{pdklib/layers\_writer.py}, verified
against the real, 377-layer \texttt{sg13g2.lyp}. Only the fields that
are genuinely real, patchable properties (name, GDS layer/datatype,
frame/fill color) are ever written back; everything else is
project-only metadata that is never round-tripped into the real file.

\subsection{Magic technology files}
\label{sec:magic}

\texttt{pdklib/magic\_tech.py} parses Magic's \texttt{.tech} format,
whose sections range from cleanly tabular
(\texttt{tech}/\texttt{version}/\texttt{planes}/\texttt{types}/\allowbreak
\texttt{contact}/\texttt{aliases}/\texttt{styles}/\texttt{compose}/\texttt{connect},
parsed in full) to small mini-rule-languages of their own
(\texttt{cifoutput}, \texttt{cifinput}, \texttt{drc}, \texttt{extract}).
For each of the latter, one or more reliable, uniformly-shaped
statement kinds are extracted in full rather than attempting a
general-purpose grammar for the whole section:

\begin{itemize}[leftmargin=1.4em]
  \item \textbf{cifoutput}: every real \texttt{layer NAME \ldots calma L D}
    recipe -- a genuine Magic-type-name to GDS-layer/datatype mapping,
    cross-referenced against the KLayout \texttt{.lyp} layers in
    \texttt{pdklib/reconcile.py} (100/377 real GDS layer/datatype pairs
    recognized by both independent descriptions of the same
    fabricated stream; the remaining 277 \texttt{.lyp}-only pairs are
    annotation/QA datatypes Magic's own fabrication output has no
    reason to paint -- a sensible finding, not a parsing gap).
  \item \textbf{cifinput}: standalone \texttt{ignore}/\texttt{calma}
    hints (23 and 133 real entries respectively), plus the section's
    own real geometry-boolean recipe blocks
    (\texttt{layer}/\texttt{templayer NAME <base> \ldots} followed by
    real \texttt{and}/\texttt{and-not}/\texttt{or}/\texttt{grow}/\allowbreak
    \texttt{shrink}/\texttt{labels}/\ldots op lines), recorded
    verbatim rather than evaluated. A real structural fact was
    confirmed here: one Magic type's own recipe can be split across
    multiple, textually distant blocks sharing the same name (e.g.\
    the real \texttt{nwell} type has two separate
    \texttt{layer nwell \ldots} blocks with two different base
    layers), mirroring a pattern already handled for
    \texttt{cifoutput} -- 188 real, unique recipe names were found
    across 211 real block occurrences once merged correctly.
  \item \textbf{drc}: the dominant tabular kinds --
    \texttt{width}/\texttt{spacing}/\texttt{maxwidth} (192 real
    statements) and \texttt{angles} (17 real statements, kept in a
    separate dataclass since a degree count is not a length and must
    not receive the same micron-conversion factor). Real values are
    converted to microns via an empirically confirmed
    \texttt{/1000} factor, cross-checked against three independently
    extracted KLayout DRC rules sharing the same rule ID.
  \item \textbf{extract}: per-layer sheet resistance and plane
    ordering, the section's own parasitic-capacitance coefficients
    (506 real lines), every real \texttt{device} statement (50 real
    entries), and its \texttt{contact}/\texttt{devresist}/\texttt{antenna}/\allowbreak
    \texttt{disconnect}/\texttt{substrate} statements (41 real lines
    combined) -- small and heterogeneous enough to share one generic,
    raw statement dataclass rather than five near-identical ones.
\end{itemize}

\subsection{LEF}

\texttt{pdklib/lef.py} parses tech-LEF layers, \texttt{SITE}s, and real
\texttt{MACRO}s (class, size, site, symmetry) with their
\texttt{PIN}s (direction, use, port layers and rectangle counts) and
\texttt{OBS} layers, hand-verified against all 32 of IHP's real
downloaded \texttt{.lef} files (a tech LEF, two combined-family LEFs,
and 28 real SRAM hard-macro LEFs). Real \texttt{VIA}/\texttt{ViaRULE}
via-stack geometry is also parsed and shown read-only. Pin data is
fully editable in both the LEF tab and the By Cell tab, through the
shared object cache described above, with a verified, byte-identical
write-back path (\texttt{pdklib/lef\_writer.py}).

\subsection{KLayout DRC decks}

\texttt{pdklib/drc.py} extracts two reliable patterns from the real
DRC-deck Ruby DSL --
\texttt{width()}/\texttt{space()}/\texttt{sep()} and
\texttt{.enclosed()} results assigned to a variable and later fed to
a same-variable \texttt{.output()} call, cross-referenced against a
real JSON values file -- yielding 75 real rules. The remaining 86
real constructs are honestly reported as skipped rather than
guessed: they are multi-step composite or derived checks (angle and
acute-corner checks, antenna-ratio accumulation, density windows)
with no single reliable, generic extraction pattern, confirmed by
tallying every real method-chain shape used across them deck-wide and
finding no one shape dominating the way the four extracted patterns
do.

\subsection{Liberty timing libraries}

\texttt{pdklib/liberty.py} extracts real per-pin and per-timing-arc
scalar data (direction, capacitance, function per pin;
related\_pin/timing\_type/timing\_sense/when per arc) across all 97
real \texttt{.lib} files and 700 real cells, editable through a
two-level dialog (Pins, then Timing Arcs for the selected pin). The
real, nested lookup-table sub-groups
(\texttt{cell\_rise}/\texttt{cell\_fall}/\texttt{rise\_transition}/\allowbreak
\texttt{fall\_transition}, each carrying its own multi-dimensional
\texttt{index\_1}/\texttt{index\_2}/\texttt{values} data) are
deliberately left unparsed, the same bounded-scope precedent used
throughout the project: a domain is either fully, verifiably
supported, or honestly reported as not yet attempted -- never
partially guessed.

\subsection{GDS}

\texttt{pdklib/gds.py} provides real structural summaries -- bounding
box and per-layer shape counts, not full polygon geometry -- via
KLayout's own Python API (\texttt{klayout.db}), lazily imported so
every other command in the toolkit still runs without it installed.
Against \texttt{sg13g2\_stdcell.gds}, 84 real structures are found,
matching LEF's own 84 macros exactly; the largest real file
(\texttt{sg13g2\_io.gds}, 69\,MB) parses in under 0.1\,s.

\subsection{Netlists, digital Verilog, and ngspice/xschem}

\texttt{pdklib/netlist.py} and \texttt{pdklib/verilog.py} extract real CDL/SPICE
subcircuits and Verilog modules with real per-port direction (from a
real CDL \texttt{*.PININFO} comment or real Verilog
\texttt{input}/\texttt{output}/\texttt{inout} declarations), editable
via a shared port editor. \texttt{pdklib/spice\_models.py} extracts real
ngspice \texttt{.model}/\texttt{.subckt} statements and, in a later
extension, real \texttt{.LIB NAME \ldots .ENDL} PVT-corner blocks
(their own \texttt{.param} overrides and \texttt{.include}
statements, kept distinct from unrelated \texttt{.param} lines living
inside a \texttt{.subckt} body). \texttt{pdklib/xschem.py} extracts a
real symbol's device-attribute block and real pin list from
xschem's \texttt{.sym} format, via a quote-aware, brace-depth-tracking
scanner needed because real content can contain escaped braces inside
quoted Tcl expressions -- 374 real pins were found across 202 real
symbols.

\section{Write-back and verification methodology}
\label{sec:verification}

Every currently structured-editable domain -- LEF pins, DRC rules,
Magic types, CDL/SPICE/Verilog ports, Liberty pin/timing data, and
KLayout layers -- has a native write-back path that is verified two
ways: first, that exporting with \emph{no edits} reproduces the real
original file byte-for-byte, across every real file in the domain
(not a sample); second, that a real, driven edit round-trips
correctly (survives export, re-import, and a simulated application
relaunch).

\subsection{Netlist/Verilog write-back}
\label{sec:netlist-writeback}

Unlike the other writers, CDL/SPICE/Verilog port write-back
(\texttt{pdklib/netlist\_writer.py}, \texttt{pdklib/verilog\_writer.py}) has
no single stable per-entry line range to patch in place: a real
\texttt{.SUBCKT} header can wrap across \texttt{+}-continuations in
ways that are unsafe to token-patch. Both writers instead compare
each cell or module's \emph{current} port list against a fresh
re-parse of the same real original file -- unchanged, the original
header text is kept byte-for-byte, wrapping included; changed, it is
regenerated cleanly as a single line, with, for CDL, a fresh
\texttt{*.PININFO} line inserted where none existed before.

\subsection{Whole-PDK, two-generation round-trip export}

Beyond per-domain verification, the toolkit's \texttt{export-full}
command exports a complete, standalone copy of the entire real PDK
tree, and this is chained twice -- export the real original to
generation A, then export A to generation B -- with a structural,
symlink-aware \texttt{diff -rq --no-dereference} comparison across
all three trees (original vs.\ A, A vs.\ B, original vs.\ B). This
caught two real bugs invisible to narrower, per-file tests: a naive
directory-copy step that silently dereferenced a real symlink
(\texttt{libs.tech/ngspice/install.py -> ../xschem/install.py}) into
a plain file copy, and a shared newline-handling bug in which every
writer unconditionally appended a trailing newline, silently
changing real byte content in the 21 real files that do not end with
one -- invisible to each writer's own narrower test because that
test's own comparison logic happened to normalize away the exact
same discrepancy it should have caught. Both were found, root-caused
against the real data, and fixed once, in one shared location, rather
than patched separately per writer.

As of this report, a full two-generation round-trip across the
complete real IHP SG13G2 tree -- 5{,}345 real files and symlinks,
744\,MB -- produces \emph{zero} differences in any of the three
pairwise comparisons.

\section{User-defined behavioral models}

Beyond editing the PDK's own native content, the toolkit supports
attaching user-authored Verilog or Verilog-A behavioral models to
real (or entirely new) cells, so a designer's own compact models can
be represented in the same cell-centric view as the foundry's native
data. Verilog-A's real \texttt{module NAME(port, \ldots);\ \ldots\ endmodule}
header and port syntax is identical to plain digital Verilog's
(confirmed against IHP's own real \texttt{mosvar.va}), so the
existing Verilog module parser is reused directly for both file
types, with no new grammar required.

Two linking mechanisms are supported: automatic name-matching (a
module named identically to a real cell attaches to it with no
configuration, the same convention every other native view already
uses for cell aggregation), and an explicit, persisted link record
mapping a specific module in a specific file to a specific cell name
-- including a wholly new, user-defined cell name not present in the
real PDK at all, which still receives its own aggregated cell-view
entry. This lets a designer author and browse an entirely new
compact model alongside the foundry's own cells, without first having
to fabricate a placeholder native view for it.

\section{Discussion and future work}

The toolkit's current scope is deliberately bounded: it reads,
cross-references, and edits an \emph{existing} real PDK's structured
content, with verified round-trip fidelity, but does not yet generate
new PDK content from scratch (a novel technology file, a new DRC
deck, a new device model) -- the stated longer-term goal. Concretely
still open: the much larger, structurally irregular remainder of
Magic's own \texttt{drc} section (variable argument counts, nested
layer-boolean expressions, and one genuinely non-check, conditional-scoping
construct); Liberty's own nested lookup-table sub-groups; the
harder, composite/derived KLayout DRC constructs with no single
reliable extraction pattern; and Qucs-S's own real schematic-symbol
format, which uses a genuinely different, unrelated tag syntax from
xschem's and would need its own, separate parser rather than a reuse
of the existing one.

\section{Conclusion}

openPDKcreator demonstrates that a general, cross-format PDK toolkit
is practical to build and verify against real, complete, open-source
foundry data -- not merely against synthetic fixtures -- and that
doing so surfaces real structural facts and real bugs that narrower,
per-format tooling and per-file testing would not have caught. Every
currently editable domain has a write-back path verified
byte-identical against the real, complete IHP SG13G2 PDK, both
per-file and at whole-tree, two-generation round-trip scale, with
zero structural drift found across 5{,}345 real files.

\end{document}
